%==============================================================================
%  13-people-skills-culture.tex — People, Skills & Culture
%==============================================================================
\strategyPart{13 \textbar\ People \& Culture}{People, Skills and Organizational Culture}

The workforce is the adoption bottleneck. This part defines AI literacy,
specialist skills, leadership capabilities and the cultural conditions that
determine whether employees use the tools.

\section{AI literacy for every employee}

\begin{itemize}
  \item Basic AI concepts and appropriate use.
  \item Data privacy and security.
  \item Prompting, verification and \textbf{hallucination awareness}.
  \item Responsible-use norms [WEF Future of Jobs; OECD; NIST AI RMF].
\end{itemize}

\section{Specialist AI skills in demand}

\begin{itemize}
  \item Data scientists and ML engineers.
  \item AI architects and AI product managers.
  \item AI governance and security specialists.
  \item AI interaction designers and evaluation engineers.
  \item Agent-orchestration engineers [WEF Future of Jobs; Stanford AI Index].
\end{itemize}

\section{Workforce productivity and skill-gap tracking}

\begin{center}
  \renewcommand{\arraystretch}{1.45}
  \rowcolors{2}{inSnow}{white}
  \fitwidth{\begin{tabular}{>{\raggedright\arraybackslash}p{3.4cm}>{\raggedright\arraybackslash}p{4.6cm}>{\raggedright\arraybackslash}p{3.6cm}}
    \rowcolor{inNavy}
    \hcell{Skill dimension} & \hcell{Current state (typical)} & \hcell{Target state} \\
    \textbf{AI literacy} & Under 30\% of employees trained & 100\% of knowledge workers \\
    \textbf{Prompt \& verification} & Ad-hoc, ungoverned & Certified, supervised \\
    \textbf{Data skills} & Data team only & Distributed data literacy \\
    \textbf{Leadership AI fluency} & Low board-level fluency & Executives able to interrogate AI \\
    \textbf{Agent supervision} & None & Defined supervision roles \\
  \end{tabular}}
\end{center}

\section{Employee adoption factors}

Adoption follows the technology-acceptance literature: \textbf{perceived
usefulness}, \textbf{ease of use}, \textbf{trust}, \textbf{training},
\textbf{managerial support}, \textbf{job-security concerns} and \textbf{quality
of output}. The organizational evidence is sobering: while adoption is
rising — 65\% of organizations now regularly use generative AI [McKinsey
2024] — value is not following automatically, because adoption without
redesigned work and incentives produces usage, not outcomes.

\section{Cultural conditions for transformation}

\begin{itemize}
  \item \textbf{Experimentation} and learning from failure.
  \item \textbf{Data-driven decisions} with transparent metrics.
  \item \textbf{Cross-functional collaboration} over silos.
  \item \textbf{Augmentation, not substitution} — the message that protects
        adoption and reduces resistance [HBR 2018].
  \item \textbf{Responsible risk-taking} with governance guardrails.
\end{itemize}
